\subsection{Definition of ordered monoids and groups}

DEFINITION 1. --- A commutative monoid structure (written additively) and an ordering (written ≤) on a set M are said to be compatible if they satisfy the following axiom:

(OM) For each \(z \in \mathbf{M}\) , the relation \(x \leqslant y\) implies \(x + z \leqslant y + z\) .

A set M equipped with a commutative monoid structure and an ordering which are compatible is called an ordered monoid; if its commutative monoid structure is a group structure, then it is called an ordered group.

In an analogous fashion, we can define the notion of a noncommutative ordered monoid (VI, p.~30, Ex. 1).

If an ordering is compatible with a given monoid structure, then so is the opposite ordering.

Examples. --- 1) The additive group of the rational integers, and that of the rational numbers, are ordered groups under the ordering defined in I, pp.~21 and 117.

* The same is true for the additive group of the real numbers (Gen.~Top., IV, p.~3). *

\begin{enumerate}
\def\labelenumi{\arabic{enumi})}
\setcounter{enumi}{1}
\tightlist
\item
  * The additive group of finite numerical functions defined on a set E is an ordered group under the ordering « \(f(x) \leq g(x)\) for all \(x \in E\) », written \(f \leq g\) . This relation says that the graph of the function f lies below that of the function g. The reader may find it convenient to think in terms of this graphical interpretation occasionally. *
\end{enumerate}

According to general definitions (Set Theory, IV, p.~264), a bijective map \(f\) from an ordered monoid \(M\) onto an ordered monoid \(M'\) is called an isomorphism of \(M\) onto \(M'\) if the structure of \(M'\) is obtained from that of \(M\) by transporting it along \(f\) . This is equivalent to saying that \(f\) is a mapping onto \(M'\) such that

\[
f (x + y) = f (x) + f (y)
\]

(that is to say a homomorphism of the monoid M onto the monoid \(\mathbf{M}'\) ), and that the relations \(x \leqslant y\) and \(f(x) \leqslant f(y)\) are equivalent (whence in particular \(f(x) = f(y)\) implies \(x = y\) , that is \(f\) is injective).

PROPOSITION 1 (« addition of inequalities »). --- In an ordered monoid M, let \((x_{i})\) and \((y_{i})\) ( \(1 \leqslant i \leqslant n\) ) be two sequences of n elements such that \(x_{i} \leqslant y_{i}\) for all i; then one has

\[
x _ {1} + \dots + x _ {n} \leqslant y _ {1} + \dots + y _ {n}.
\]

If moreover all the elements \(x_{i}, y_{i}\) are cancellable (I, p.~15, Def. 5) (in particular if \(M\) is a group), and if \(x_{i} < y_{i}\) for some \(i\) , then

\[
x _ {1} + \dots + x _ {n} <   y _ {1} + \dots + y _ {n}.
\]

The general case reduces by induction to the case n = 2, using the fact that a sum of cancellable elements is cancellable for the second assertion (I, p.~15, Prop. 2). The first assertion follows from the relations

\[
x _ {1} + x _ {2} \leqslant x _ {1} + y _ {2} \quad \text { and } \quad x _ {1} + y _ {2} \leqslant y _ {1} + y _ {2},
\]

which are consequences of the hypotheses and of (OM). This being the case, the relation

\[
x _ {1} + x _ {2} = y _ {1} + y _ {2}
\]

would imply

\[
x _ {1} + x _ {2} = x _ {1} + y _ {2} = y _ {1} + y _ {2},
\]

whence \(x_{2} = y_{2}\) and \(x_{1} = y_{1}\) if \(x_{1}\) and \(y_{2}\) are cancellable, which proves the second assertion.

PROPOSITION 2. --- In an ordered group G the relations \(x \leqslant y\) and \(x + z \leqslant y + z\) are equivalent.

Indeed one can obtain each from the other by adding z or \((-z)\) to both sides.

This fact can be expressed by saying that, in an ordered group G, the ordering is translation-invariant. In other words a translation is an automorphism of the ordering of an ordered group.

COROLLARY. --- In an ordered group G, the relations \(x \leqslant y\) , \(0 \leqslant y - x\) , \(x - y \leqslant 0\) and \(-y \leqslant -x\) are equivalent.

In fact, we apply Prop. 2, successively taking \(z = -x\) , \(z = -y\) and \(z = - (x + y)\) .

In particular we deduce from this corollary that if G is an ordered group, the map \(x \mapsto -x\) from G to itself takes the ordering of G to the opposite ordering.

